\documentclass[11pt,a4paper]{article}
\usepackage[margin=25mm]{geometry}
\usepackage{amsmath,amssymb,amsthm,hyperref}
\setlength{\emergencystretch}{2em}
\newtheorem{theorem}{Theorem}
\newtheorem{corollary}[theorem]{Corollary}
\newcommand{\gen}[1]{\langle#1\rangle}
\newcommand{\Alt}{\operatorname{Alt}}
\title{A pure head cannot attach trivially to a natural alternating tail}
\author{Human Compression project: L044 continuation}
\date{}
\begin{document}\maketitle
Let $\Omega$ be finite and $G=\gen{\rho_0,\ldots,\rho_{r-1}}\le\Alt(\Omega)$ be a
string C-group. Suppose $\Omega=D\sqcup E$, $|D|\ge2$, $|E|\ge3$,
and for some $i$ its actual supported subgroups satisfy
\[
 A=G_{<i}\text{ is transitive on }D\text{ and fixes }E,
 \qquad B=G_{>i}=\Alt(E)\text{ fixes }D.
\]
Put $\tau=\rho_i$, $\sigma=\rho_{i+1}$ and $H=G_{>i+1}$.
Assume that for $v\in D$ and $b\in E$,
\[
          \tau=(v\ b)\delta,\qquad
          \operatorname{supp}(\delta)\subseteq E\setminus\{b\}.
                                                               \tag{1}
\]
In particular the crossing generator fixes every other point of the
head. The requirement (1) is essential; a nontrivial head-side factor
of the crossing generator is not silently dropped.

\begin{theorem}
No such even string C-group exists.
\end{theorem}
\begin{proof}
The group $H$ fixes $D$ and centralizes $\tau$, so fixes $b$.
Transitivity of $B=\gen{\sigma,H}$ gives $x=b^\sigma\ne b$;
$\sigma$ fixes $D$, so $x\in E\setminus\{b\}$.
The involution $\delta$ is odd, since $\tau$ is even.

For each $u\in D$, transitivity of $A$ supplies the conjugate
$(u\ b)\delta$ of $\tau$. Products of two such conjugates are the
pure three-cycles on $\{b,u,w\}$, for distinct $u,w\in D$.
They generate the natural $\Alt(D\cup\{b\})$, fixing its complement.
Conjugate by $\sigma$. These two natural alternating groups have
common set $D$ of size at least two, so they generate the natural
alternating group on their union. Hence
\[
 \Alt(D\cup\{b,x\})\le P=G_{\le i+1}.                    \tag{2}
\]
The alternating-overlap fact used here is elementary: the pure
three-cycles on the two supports, conjugated using two common
points, supply the three-cycles across the supports. It is valid
also when either initial support has three points.

Put $Z=E\cup\{v\}$. Both $B$ and $B^\tau$ are natural alternating
groups on subsets of $Z$ of size $|E|$ with $|E|-1\ge2$ common
points. Their union is $Z$. Since $\tau$ is even and fixes
$D\setminus\{v\}$, the same overlap argument gives
\[
                       R=G_{\ge i}=\Alt(Z).             \tag{3}
\]
Actual IP therefore gives $P\cap R=\gen{\tau,\sigma}$.

We first prove that $\tau$ fixes $x$. If $y=x^\tau\ne x$,
then $y\notin D\cup\{b,x\}$, because $\tau$ preserves
$D\cup\{b\}$. Conjugate (2) by $\tau$ and use overlap once
more. This puts pure $\Alt(D\cup\{b,x,y\})$ in $P$.
All four points $v,b,x,y$ lie in $Z$, so (3) and full IP put
pure $\Alt(\{v,b,x,y\})$ in $\gen{\tau,\sigma}$.
A group generated by two involutions is dihedral, and all its
odd-order elements lie in its cyclic rotation subgroup. But the
three-cycles on $\{v,b,x\}$ and $\{v,b,y\}$ do not commute.
This contradiction proves $x^\tau=x$, and hence $x^\delta=x$.
This is the complete already used pure-head fixation argument.

By (2)--(3), the pure three-cycle $q=(v\ b\ x)$, or its inverse,
belongs to $\gen{\tau,\sigma}$. Choose its orientation so that
\[
                          \tau^q=(b\ x)\delta.
\]
Write $\sigma=(b\ x)\epsilon$, where $\epsilon$ fixes
$D\cup\{b,x\}$ pointwise. Both $\delta$ and $\epsilon$ are odd.
Their product
\[
 w=\tau^q\sigma=\delta\epsilon
\]
is even and supported on $E$, and hence belongs to the actual
$B=\Alt(E)$. It also belongs to $\gen{\tau,\sigma}$.
Full IP now gives
\[
          w\in\gen{\tau,\sigma}\cap B=\gen{\sigma}.
\]
Since $w$ fixes $b$ and $\sigma$ moves it, $w=1$.
Thus $\epsilon=\delta$ and $\sigma=(b\ x)\delta$.

The involution $\delta$ centralizes $H$: the latter fixes $v,b$
and commutes with $\tau=(v\ b)\delta$. It also centralizes
$\sigma=(b\ x)\delta$, because it fixes $b,x$.
Consequently it centralizes the transitive group
$B=\gen{\sigma,H}$ on $E$, while fixing $b\in E$.
A centralizing permutation of a transitive action that fixes one
point fixes every point. Hence $\delta=1$, contradicting its
oddness. This proves the theorem.
\end{proof}

\begin{corollary}[The actual alternating-side subcase of source 5.31--5.32]
Retain the source's supported-interval identification
$A=G_{<i}$ and $B=G_{>i}$, acting on the two omission orbits
$D,E$, respectively. Suppose $A$ is the actual six-point head of
Table~2(2) or (3), allowing full reversal of its three labels,
and $B$ is the natural alternating group on $E$, with $|E|\ge3$.
Suppose the crossing generator has exactly one edge between $D$ and $E$.
Then this configuration cannot occur in an even string C-group.
The fully reversed ambient configuration is excluded as well.
\end{corollary}
\begin{proof}
Here $i=3$. In the displayed orientation the first two head
involutions are $(1\ 2)(3\ 4)$ and $(2\ 3)(4\ 5)$.
They generate $D_{10}$ on $\{1,\ldots,5\}$, with trivial
centralizer on that set, and fix point 6. Their unique nontrivial
orbit is preserved by $\rho_3$, which commutes with both.
Thus $\rho_3$ fixes points 1 through 5. The unique crossing has
endpoint 6, so $\rho_3=(6\ b)\delta$ with
$\operatorname{supp}(\delta)\subseteq E\setminus\{b\}$.
All theorem hypotheses now hold.

For clarity the reversed head is not silently identified with
its original ordered tuple. In reversed row 2, the first two
involutions have two three-point orbits covering $D$, so a commuting
$\rho_3$ preserves $D$: all its other orbits are singletons and
cannot be exchanged with these triples. It cannot have the assumed
crossing. Reversed row 3 is conjugate to its displayed orientation
by $(1\ 6)(2\ 5)(3\ 4)$, giving the same argument.
These literal head facts are retained from the accepted ordered
Table~2 and critical-head proofs; no new enumeration is used.
\end{proof}

\paragraph{Source interface and attribution.}
Cameron--Fernandes--Leemans--Mixer,
\emph{Highest rank of a polytope for $A_n$}, arXiv:1605.09173v1,
physical page 26 explicitly identifies the supported sides with
$G_{<i}$ and $G_{>i}$ in Section 5.2. The exceptional cases of
that section are sent to Case 3 on page 29. Propositions 5.31
(page 38) and 5.32 (page 44) explicitly assume the opposite side
is alternating. The corollary covers their six-point-head subcase
with this actual supported identification. It does not equate
an arbitrary projected alternating side with the whole actual suffix.

The accepted general critical-head reduction supplies
$R=G_{\ge3}=\Alt(Z)$, but does \emph{not} by itself supply
$B=G_{\ge4}=\Alt(E)$. Thus that reduction alone is insufficient
for this application. The stronger $B$ assumption is explicit here;
(3) then derives $R$ independently. Neither group is identified
merely from an abstract isomorphism or an orbit projection.
The first pure-alternating and dihedral step is reproduced from the
accepted strict-internal head lemma. The added three-cycle conjugation
and full-IP extraction force the central residual involution to vanish.
The theorem needs neither a rank estimate nor a finite census.
Arbitrary exceptional ends without the natural $B$ premise,
remaining singleton routing, connected-square rigidity and table
contracts are separate. No source-wide, formalization, question or
historical-priority claim is made.
\end{document}
